\section{Introduction}
An experiment can admit a small positive factorization at each separate
cut while requiring substantially more states when the factorizations must
be joined by common stochastic updates. For hidden realizations, a state
need not carry an individually meaningful predictive vector. A projection
of a $K$-coordinate simplex section can have exponentially many vertices.
Consequently, a lower bound obtained only by counting the vertices of an
observable polygon may lose an exponential factor in the memory parameter.

We study the following fixed-alphabet experiment. A seed
$x\in\{-1,1\}^2$ is followed by $N$ commands $c_t\in\{0,1\}$,
a query $j\in\{1,2\}$, and a binary answer $B$. For
$\zeta=e^{2\pi i\alpha}$ and the corresponding planar rotation $R_\alpha$,
the prescribed answer law is
\begin{equation}\label{eq:array34}
 \Pp(B=b\mid x,c_1,\ldots,c_N,j)
   =\frac{1+b\,e_j^TR_\alpha^{c_1+\cdots+c_N}x/10}{2}.
\end{equation}
All these probabilities lie in the fixed interval
$[(1-\sqrt2/10)/2,(1+\sqrt2/10)/2]$. The alphabets, amplitude and
support margin do not change with $N$. Exact realization means equality
for every input word, not merely after averaging over commands.

The separate positive minimum of this experiment is three. The previous
mean-support argument gave $\Omega_\alpha(\log N)$ as an unrestricted
hidden-state lower bound and $\Theta_\alpha(N^{1/3})$ within a vector-state
subclass at bounded-type angles \cite{GTF33}. We retain its exact upper
construction but replace the hidden lower argument. The new method does
not associate a predictive vector with each basis state. Instead, it
studies the conditional expectation of a random past phase \emph{under a
chosen test distribution on external words}. These conditional expectations
are proof objects, not information made available to the executing machine.

Write
\begin{equation}\label{eq:beta34}
 \beta_m(\alpha)=\min_{1\le l\le m}\sin^2(\pi l\alpha),\qquad
 b_\alpha=\inf_{l\ge1}l\|l\alpha\|_{\R/\mathbb Z}.
\end{equation}
For irrational $\alpha$, every $\beta_m$ is positive. The assumption
$b_\alpha>0$ is the usual badly-approximable, equivalently bounded-type,
condition. Let $W_{N,\varepsilon}(\alpha)$ be the least peak label count
of a clocked stochastic realization with uniform row-total-variation error
at most $\varepsilon$, and write $W_N=W_{N,0}$.

\begin{theorem}[Polynomial obstruction for all hidden realizations]
\label{thm:main34}
Let $0\le\varepsilon<1/20$ and put $\kappa=1/10-2\varepsilon$.
Every realizing profile $(K_t)_{t=0}^N$ satisfies, for every integer $k\ge1$,
\begin{equation}\label{eq:main-count34}
 \#\{0\le t<N:K_t\le k\}\,\beta_{2k}(\alpha)
                       \le \frac{18k^3}{\kappa^2}.
\end{equation}
In particular, if $b_\alpha>0$, then
\begin{equation}\label{eq:main-lower34}
 W_{N,\varepsilon}(\alpha)
       \ge \left(\frac{b_\alpha^2\kappa^2 N}{18}\right)^{1/5}.
\end{equation}
For $\alpha=(\sqrt5-1)/2$, exact realization therefore obeys
\begin{equation}\label{eq:golden34}
 \max\{3,(N/16200)^{1/5}\}\le W_N(\alpha)
       <2\max\{5,(40N)^{1/3}\}.
\end{equation}
More generally the same cube-root upper order holds whenever the partial
quotients of $\alpha$ are bounded. For $\zeta=(3+4i)/5$,
\begin{equation}\label{eq:rational-head34}
 N\le7200 W_N^3 25^{2W_N},\qquad
 W_N\ge\max\{3,\log_{5000}(N/7200)\}.
\end{equation}
All lower bounds allow arbitrary cut-dependent hidden-state continuations.
\end{theorem}

The theorem excludes $O(\log N)$ hidden width at bounded-type angles,
including at any fixed row error below $1/20$. It does not prove that
hidden and vector-state widths have the same order: the current exact
bounds are $\Omega_\alpha(N^{1/5})$ and $O_\alpha(N^{1/3})$.
The cumulative version \eqref{eq:main-count34} also limits how often a
small register may occur in an otherwise unrestricted profile.

The proof has three steps. Conditional phase amplitude decreases under a
stochastic update. Its entire loss is at most one. A homogeneous angular
inequality bounds the perturbation of the $l$th Fourier moment by
$2l^2$ times this same loss. Finally, the positive Fej\'er kernel forces
a nonnegligible Fourier potential whenever the conditional phase measure
has few atoms. Randomizing the two external commands damps those moments
at the rates $\sin^2(\pi l\alpha)$. The resulting potential telescopes
for time-dependent as well as stationary matrices.

Section~\ref{sec:deficiency34} interprets a supplied lift diagram as a
sequence of finite experiment comparisons. The one-step garbling criterion
and deficiency dual are classical \cite{Blackwell,Torgersen}. Their
composition produces an executable approximate transducer. Applying
Theorem~\ref{thm:main34} yields a new lower bound on the total defect of
\emph{every} such diagram for \eqref{eq:array34}, rather than only a
four-point warning example. Section~\ref{sec:symmetry34} compares this
notion with equivariant polyhedral lifts and distinguishes fixed-length
autonomy from the stronger anytime query task.

\begin{center}\small
\begin{tabular}{p{.19\textwidth}p{.72\textwidth}}
\toprule
Notation & Meaning \\
\midrule
$K_t$, $W_{N,\varepsilon}$ & Available labels at cut $t$; minimum peak at row error $\varepsilon$.\\
$z_t(s)$, $q_t$ & Conditional past-phase expectation; its mean absolute value.\\
$\Phi_t^{(m)}$ & Finite Fej\'er-weighted Fourier potential of the conditional phases.\\
$\delta(E,F)$ & One-sided finite deficiency: simulate $F$ by garbling $E$.\\
$\mathcal D(\mathcal E)$ & Sum of optimized one-step and terminal defects for a fixed diagram.\\
$A_N^{\rm fix}$, $A_N^{\rm any}$ & Autonomous state counts for one fixed length, or all lengths up to $N$.\\
\bottomrule
\end{tabular}
\end{center}
